% =========================================================================
% ภาคผนวก ฉ: อภิธานศัพท์และคำย่อเฉพาะทาง 2 ภาษา (Bilingual Glossary & Acronyms)
% Appendix F: Bilingual Glossary of Technical Terms and List of Acronyms
% =========================================================================

\chapter{อภิธานศัพท์และคำย่อเฉพาะทาง 2 ภาษา (Bilingual Glossary \& Acronyms)}
\label{app:glossary_and_acronyms}

ภาคผนวกนี้รวบรวมคำย่อสากล ศัพท์บัญญัติ และคำนิยามเชิงเทคนิคที่ใช้ในตำราเล่มนี้ เพื่ออำนวยความสะดวกในการค้นคว้า ทำความเข้าใจ และอ้างอิงตามมาตรฐานวิชาการของราชบัณฑิตยสภาและองค์กรวิทยาศาสตร์สากล

\section{รายการคำย่อสากล (List of Acronyms)}
\label{app:sec_acronyms}

\begin{table}[H]
\centering
\caption{รายการคำย่อสากลทางฟิสิกส์พลาสมาและปัญญาประดิษฐ์ (ตอนที่ 1: A--L)}
\label{tab:app_acronyms_part1}
\begin{tabularx}{\textwidth}{l l X}
\toprule
\textbf{คำย่อ} & \textbf{คำเต็มภาษาอังกฤษ} & \textbf{ความหมายภาษาไทย} \\
\midrule
\textbf{ACGIH} & American Conf. of Govt. Ind. Hygienists & สมาคมนักสุขศาสตร์อุตสาหกรรมภาครัฐสหรัฐฯ \\
\textbf{CAP} & Cold Atmospheric Plasma & พลาสมาความดันบรรยากาศอุณหภูมิต่ำ (พลาสมาเย็น) \\
\textbf{CEC} & Cation Exchange Capacity & ความจุในการแลกเปลี่ยนประจุบวกของดิน \\
\textbf{CPPS} & Cyber-Physical Plasma System & ระบบพลาสมากายภาพไซเบอร์ \\
\textbf{DBD} & Dielectric Barrier Discharge & การคายประจุผ่านแนวกั้นไดอิเล็กทริก \\
\textbf{EC} & Electrical Conductivity & สภาพนำไฟฟ้าของสารละลาย \\
\textbf{EEDF} & Electron Energy Distribution Function & ฟังก์ชันการแจกแจงพลังงานของอิเล็กตรอน \\
\textbf{GaN} & Gallium Nitride & แกลเลียมไนไตรด์ (สารกึ่งตัวนำแถบพลังงานกว้าง) \\
\textbf{GC-MS} & Gas Chromatography--Mass Spectrometry & แก๊สโครมาโทกราฟี-แมสสเปกโทรเมตรี \\
\textbf{HPLC} & High-Performance Liquid Chromatography & โครมาโทกราฟีของเหลวสมรรถนะสูง \\
\textbf{LCA} & Life Cycle Assessment & การประเมินวัฏจักรชีวิต \\
\textbf{LCOE} & Levelized Cost of Energy & ต้นทุนพลังงานปรับระดับ \\
\bottomrule
\end{tabularx}
\end{table}

\begin{table}[H]
\centering
\caption{รายการคำย่อสากลทางฟิสิกส์พลาสมาและปัญญาประดิษฐ์ (ตอนที่ 2: N--V)}
\label{tab:app_acronyms_part2}
\begin{tabularx}{\textwidth}{l l X}
\toprule
\textbf{คำย่อ} & \textbf{คำเต็มภาษาอังกฤษ} & \textbf{ความหมายภาษาไทย} \\
\midrule
\textbf{NIOSH} & National Inst. for Occupational Safety & สถาบันแห่งชาติด้านความปลอดภัยและอาชีวอนามัย \\
\textbf{OES} & Optical Emission Spectroscopy & สเปกโทรสโกปีการปลดปล่อยเชิงแสง \\
\textbf{ORP} & Oxidation-Reduction Potential & ศักย์ออกซิเดชัน-รีดักชัน \\
\textbf{OSHA} & Occupational Safety and Health Admin. & องค์การบริหารความปลอดภัยและอาชีวอนามัย \\
\textbf{OWRK} & Owens, Wendt, Rabel, and Kaelble & แบบจำลองคำนวณพลังงานอิสระที่พื้นผิว \\
\textbf{PAW} & Plasma-Activated Water & น้ำกระตุ้นด้วยพลาสมา \\
\textbf{PEL} & Permissible Exposure Limit & ขีดจำกัดการสัมผัสสารเคมีที่ยอมรับได้ \\
\textbf{PFAS} & Per- and Polyfluoroalkyl Substances & สารเปอร์และโพลีฟลูออโรอัลคิล \\
\textbf{PINN} & Physics-Informed Neural Network & โครงข่ายประสาทเทียมผสานฟิสิกส์ \\
\textbf{RONS} & Reactive Oxygen and Nitrogen Species & สปีชีส์ไวต่อปฏิกิริยาของออกซิเจนและไนโตรเจน \\
\textbf{SiC} & Silicon Carbide & ซิลิคอนคาร์ไบด์ (สารกึ่งตัวนำกำลังสูง) \\
\textbf{SOP} & Standard Operating Procedure & ขั้นตอนการปฏิบัติงานมาตรฐาน \\
\textbf{SVI} & Seed Vigor Index & ดัชนีความแข็งแรงของเมล็ดพันธุ์ \\
\textbf{TLV} & Threshold Limit Value & ค่าขีดจำกัดความเข้มข้นสารเคมีที่ปลอดภัย \\
\textbf{TWA} & Time-Weighted Average & ค่าเฉลี่ยถ่วงน้ำหนักตามเวลาการทำงาน 8 ชม. \\
\textbf{VOCs} & Volatile Organic Compounds & สารประกอบอินทรีย์ระเหยง่าย \\
\bottomrule
\end{tabularx}
\end{table}

\section{อภิธานศัพท์และคำนิยามเชิงเทคนิค 2 ภาษา (Bilingual Glossary)}
\label{app:sec_glossary}

\begin{description}[leftmargin=2.0em, itemsep=6pt]
    \item[\textbf{การกำบังเดอบาย (Debye Shielding):}] ปรากฏการณ์ที่อนุภาคมีประจุในพลาสมาจัดเรียงตัวใหม่เพื่อหักล้างสนามไฟฟ้าสถิตจากประจุแปลกปลอม ทำให้ศักย์ไฟฟ้าลดลงแบบเอกซ์โพเนนเชียลตามระยะทาง
    \item[\textbf{การคายประจุผ่านแนวกั้นไดอิเล็กทริก (Dielectric Barrier Discharge, DBD):}] การคายประจุไฟฟ้าในก๊าซที่มีฉนวนไฟฟ้ากั้นขวางอย่างน้อยหนึ่งระนาบ เพื่อจำกัดการไหลของกระแสไม่ให้เกิดการอาร์กความร้อนสูง
    \item[\textbf{การถล่มทลายของทาวน์เซนด์ (Townsend Avalanche):}] กลไกการทวีคูณจำนวนอิเล็กตรอนแบบลูกโซ่จากการชนแตกตัวในก๊าซภายใต้อิทธิพลของสนามไฟฟ้าสูง
    \item[\textbf{ความยาวเดอบาย (Debye Length):}] ระยะทางลักษณะเฉพาะที่ศักย์ไฟฟ้าสถิตในพลาสมาลดลงเหลือ $1/e$ ของค่าเริ่มต้น เป็นเกณฑ์บ่งชี้ความเป็นกลางทางไฟฟ้าในระดับมหภาค
    \item[\textbf{น้ำกระตุ้นด้วยพลาสมา (Plasma-Activated Water, PAW):}] สารละลายน้ำที่ผ่านการทำปฏิกิริยากับพลาสมาเย็น ทำให้มีค่า pH เป็นกรดอ่อน มีค่า ORP สูง และอุดมไปด้วยสปีชีส์ RONS เช่น $\ce{NO3^-}, \ce{NO2^-}, \ce{H2O2}$ และ $\ce{ONOOH}$
    \item[\textbf{พลาสมาเย็น (Cold Atmospheric Plasma, CAP):}] พลาสมาที่ไม่สมดุลทางความร้อน (Non-thermal plasma) ซึ่งอุณหภูมิอิเล็กตรอนมีค่าสูงมาก ($T_e \sim 1$--$5\text{ eV} \approx 10^4\text{ K}$) แต่อุณหภูมิของก๊าซและไอออนใกล้เคียงกับอุณหภูมิห้อง ($T_g \sim 300\text{ K}$)
    \item[\textbf{ฟังก์ชันการแจกแจงพลังงานอิเล็กตรอน (Electron Energy Distribution Function, EEDF):}] ฟังก์ชันความน่าจะเป็นที่อธิบายสัดส่วนของอิเล็กตรอนที่มีระดับพลังงานจลน์ต่างๆ ในพลาสมา เป็นหัวใจสำคัญในการคำนวณอัตราการเกิดปฏิกิริยาเคมี
    \item[\textbf{วงจรพัลส์มาร์กโซลิดสเตต (Solid-State Marx Generator):}] วงจรทวีคูณแรงดันไฟฟ้าที่ประจุตัวเก็บประจุแบบขนานด้วยแรงดันต่ำ แล้วสลับการคายประจุแบบอนุกรมด้วยสวิตช์สารกึ่งตัวนำกำลัง (SiC/GaN) เพื่อผลิตพัลส์ไฟฟ้าแรงดันสูงมาก
    \item[\textbf{สเปกโทรสโกปีการปลดปล่อยเชิงแสง (Optical Emission Spectroscopy, OES):}] เทคนิคการวิเคราะห์เชิงแสงแบบไม่ทำลายตัวอย่าง โดยการวัดความยาวคลื่นและความเข้มของแสงที่ปลดปล่อยออกมาจากอะตอมและโมเลกุลในพลาสมาเมื่อเปลี่ยนสถานะพลังงาน
    \item[\textbf{โครงข่ายประสาทเทียมผสานฟิสิกส์ (Physics-Informed Neural Network, PINN):}] สถาปัตยกรรมการเรียนรู้เชิงลึกที่นำสมการอนุพันธ์ย่อย (PDE) ทางฟิสิกส์ เช่น กฎการอนุรักษ์มวลและโมเมนตัม เข้าไปเป็นส่วนหนึ่งของฟังก์ชันการสูญเสียระหว่างการฝึกโมเดล
\end{description}

\clearpage
